% !TeX root = ../main.tex

\xchapter{绪论}{Introductions}

\xsection{选题意义及应用背景}{Backgrounds}

在当今数字化信息爆炸的时代，数字视频内容已经渗透到人类生活的方方面面。从传统的电视节目、电影艺术到新兴的短视频、在线直播以及游戏录像，多媒体资源以其直观、生动的特性，极大地丰富了人们的日常生活，并成为信息传播与娱乐消费的核心载体。然而，视频数据的海量增长也带来了一个严峻的挑战：信息过载。面对浩如烟海的视频资源，要求用户耗费大量时间完整观看每一个视频以获取关键信息已成为一种奢侈的负担。因此，如何在保留视频核心内容的前提下，通过自动化技术对冗长的视频进行“瘦身”，已成为计算机视觉和多媒体处理领域亟待解决的热点问题。视频摘要（Video Summarization）技术应运而生，其旨在通过算法自动提取视频中最具代表性、最显著的帧或片段，生成一段精炼的摘要，从而帮助用户在极短的时间内理解视频的主旨大意。

从技术发展的演进历程来看，传统的视频摘要方法主要依赖于监督学习、非监督学习以及弱监督学习范式。监督学习方法\cite{huang2023causalainer, zhao2017hierarchical, chen2022video}通常依赖于大量由人工标注的关键帧或视频片段作为真实标签，通过深度神经网络学习视频内容的显著性特征，并基于显著性特征完成最终的视频摘要任务。虽然这类方法在特定数据集上表现优异，但其对大规模高质量标注数据的依赖限制了模型的扩展性与迁移能力。非监督学习方法\cite{liang2022video,mahasseni2017unsupervised, apostolidis2022summarizing}则试图通过挖掘视频内部的固有属性（如帧间相似性、稀疏性或内容多样性）来生成摘要，尽管缓解了数据标注的压力，但往往难以捕捉高层语义信息。弱监督方法\cite{ye2021temporal,narasimhan2022tl}则作为一种折中方案，利用粗粒度标签或用户交互信号进行引导，在一定程度上降低了对精细标注的依赖，同时兼顾了模型训练的可行性与摘要质量。相比完全监督方法，弱监督范式具有更强的应用灵活性，但其性能上限仍然受限于弱标签信息的不完整性与噪声干扰。然而，无论是哪种传统范式，通常都需要复杂的网络设计和高昂的离线训练过程，且模型往往被限制在特定的领域，难以直接部署到计算资源有限或领域多变的实际生产环境中。

近年来，多模态大语言模型（Multimodal Large Language Models, MLLMs）的突破性进展，如 GPT-4V\cite{achiam2023gpt}、LLaVA-Video\cite{zhang2024video}等，为视觉理解任务开辟了全新的路径。这些模型通过在海量跨模态数据上的预训练，展现出了卓越的泛化能力和逻辑推理能力，使得研究者能够将视觉任务重新构建为基于提示（Prompt）的自然语言任务。目前，这种基于大模型的零样本（Zero-shot）推理范式已在视频异常检测\cite{zanella2024harnessing}、视频时序定位\cite{zheng2024training}和视频时刻检索\cite{luo2024zero}等领域取得了显著成功。受此启发，探索如何利用大语言模型的通用知识来实现无需训练（Training-free）的零样本视频摘要技术，具有重要的学术价值和实际应用前景。

本课题的研究背景正基于于这一技术变革。传统的视频摘要模型虽在特定任务上精通，但灵活性不足；而基于多模态大模型的视频摘要框架则能够凭借其深厚的语义理解基础，通过自然语言指令直接从原始视频中筛选出符合人类审美的关键片段。这种无需训练的架构具有显著优势：首先，它完全省去了耗时的模型微调过程，降低了算力门槛；其次，它具备极强的跨领域迁移能力，能够适应不同风格的视频内容；最后，它支持以语言驱动的交互方式，使得摘要的生成过程更加直观和可解释。在实际应用场景中，如新闻资讯自动生成、安防监控长视频快速回溯、个人云盘视频智能管理等领域，这种高效、低成本且插拔式的视频摘要技术将极具竞争力。因此，开发一种能够充分利用多模态大模型语义对齐能力的训练无关型视频摘要网络，不仅是响应工业界降本增效的需求，更是推动视频智能理解向通用化、交互化方向演进的关键一步。

\xsection{国内外研究现状分析}{Related Works}

视频摘要旨在通过提取视频中最具代表性的帧或片段，生成能够覆盖原始视频核心语义的紧凑表达。随着深度学习技术的发展，该领域经历了从基于手工特征的方法、到深度序列建模方法，再到当前由大语言模型驱动的多模态推理范式的演进。本节将围绕传统视频摘要方法、大语言模型在视频任务中的应用，以及无须训练的视频分析框架三个方面展开综述。

\xsubsection{早期视频摘要方法研究}{Current Status of Video Summarization Methods}
根据是否依赖人工标注数据，视频摘要方法通常分为监督学习与无监督学习两大类。早期的监督学习方法主要基于循环神经网络（Recurrent Neural Networks, RNN）\cite{rumelhart1986learning}及其变体长短期记忆网络（Long Short-Term Memory, LSTM）\cite{graves2012long}。由于视频数据天然具有时间序列特性，RNN能够有效建模帧间的时序依赖关系，因此成为初期视频摘要模型的核心架构。Zhang等人\cite{zhang2016video} 首次将LSTM与行列式点过程（Determinantal Point Process, DPP）\cite{kulesza2012determinantal}结合，在建模时间相关性的同时显式约束摘要的多样性。随后，研究者进一步探索更复杂的时序建模方式，例如通过层级化结构捕捉多尺度时间依赖\cite{zhao2017hierarchical}，以及通过结构自适应机制增强模型对不同视频内容的泛化能力\cite{zhao2018hsa}。

然而，随着视频长度和语义复杂度的增加，RNN在建模长距离依赖时的局限性逐渐显现。尽管LSTM在一定程度上缓解了梯度消失问题，但其全局建模能力仍然受限。为此，注意力机制被引入视频摘要任务，通过显式建模关键帧的重要性权重来增强语义表达能力\cite{lebron2018video,ji2020deep,fajtl2019summarizing,li2021exploring}。在此基础上，一部分工作将视频摘要建模为序列到序列（Seq2Seq）问题，利用编码器-解码器结构生成摘要序列\cite{ji2019video,rochan2018video,lal2019online}，从而提升模型在复杂场景下的表达能力。

随着Transformer架构\cite{vaswani2017attention}在自然语言处理领域取得突破，其基于自注意力机制的全局建模能力也被广泛引入视频摘要任务中。相比RNN的递归结构，Transformer可以直接建模任意两帧之间的依赖关系，从而更适合处理长视频序列。近期研究不仅利用堆叠的自注意力层进行时间建模\cite{argaw2024scaling,chen2022video,lin2023videoxum}，还进一步拓展至空间与时空联合建模，通过引入空间结构信息提升表示能力\cite{liang2022video,liu2022video,hsu2023video,son2024csta}。与此同时，为缓解Transformer计算复杂度高的问题，研究者提出了稀疏注意力\cite{child2019generating}、局部窗口机制\cite{liu2021swin}以及改进位置编码等优化策略，在性能与效率之间取得平衡。

相较于依赖大量标注数据的监督方法，无监督视频摘要旨在在无人工标注的条件下学习关键帧选择策略。该类方法的核心思想在于：一个优质摘要应能够在压缩信息的同时重建原始视频的主要内容。基于这一思想，重构驱动的生成式模型成为主流方向，典型方法包括结合变分自编码器（Variational Autoencoder, VAE）\cite{kingma2013auto}与生成对抗网络（Generative Adversarial Network, GAN）\cite{goodfellow2014generative}的统一框架\cite{mahasseni2017unsupervised,apostolidis2019stepwise}。这些方法通常利用VAE建模潜在分布，同时通过对抗训练增强摘要的代表性与多样性。进一步的研究引入注意力机制以强化依赖建模，例如通过自注意力捕捉全局与局部关系\cite{jung2020global}，或通过多模态融合减少视觉与语义之间的差距\cite{lin2022deep}。尽管无监督方法降低了对标注数据的依赖，但其性能往往依赖于复杂的训练过程，且在跨领域场景中的泛化能力仍然有限。

\xsubsection{大语言模型在视频任务中的应用研究}{Research on the Application of Large Language Models in Video Tasks}
近年来，大语言模型（Large Language Models, LLM）如GPT-4\cite{achiam2023gpt}、LLaMA\cite{touvron2023llama} 以及Qwen2\cite{yang2024qwen2}的出现，为视频理解任务提供了强大的语义建模能力。在此基础上，多模态大语言模型（Multimodal Large Language Models, MLLMs）通过将视觉信息映射到语言语义空间，实现了跨模态的统一表示，进一步推动了视频理解研究的发展。

在通用视频理解与视频对话方面，VideoChat\cite{li2023videochat}通过引入可学习的神经接口，将视频编码器与LLM进行耦合，实现了以对话为中心的视频理解框架，能够完成时空推理与事件定位等复杂任务。针对长视频带来的上下文限制问题，MA-LLM\cite{he2024ma} 提出基于外部记忆库的增强架构，通过存储历史帧特征实现长期依赖建模，同时避免上下文窗口溢出。MovieChat\cite{song2024moviechat} 则借鉴认知心理学中的记忆模型，引入稀疏记忆机制，将模型处理视频数据过程中产生的基本语义Token视为记忆单元，从而在保证效率的同时实现长时序推理。

在视频异常检测任务中，基于LLM的方法利用其强大的语义推理能力，对异常行为进行建模与解释。Holmes-VAD\cite{zhang2024holmes}通过大规模指令微调，使模型不仅能够识别异常，还能生成可解释的推理过程。AnomalyGPT\cite{gu2024anomalygpt}则通过模拟工业数据并结合提示学习，实现了无需人工阈值设定的细粒度异常检测。此外，相关研究还将LLM应用于语义异常检测，通过建模系统层面的语义偏差来识别复杂异常\cite{elhafsi2023semantic}。

在时序定位任务中，LLM同样展现出强大的能力。ChatVTG\cite{qu2024chatvtg}利用LLM生成多粒度视频片段描述，并通过语义匹配实现零样本定位，随后通过精细化步骤提升定位精度。VTimeLLM\cite{huang2024vtimellm}则提出三阶段训练策略，包括跨模态特征对齐、时间边界学习以及指令微调，从而实现对视频关键时刻的精准捕捉。这些研究表明，LLM正在将视频理解从传统的特征提取及后续结果匹配问题，转变为基于语义理解与逻辑推理的高级认知任务。

\xsubsection{免训练视频理解框架研究现状}{Research on Training-Free Video Understanding Frameworks}
尽管基于微调（Fine-tuning）或适配器（Adapter）的方法在性能上表现优越，但其训练成本和计算资源需求较高，限制了实际部署。近年来，免训练（Training-Free）的视频理解框架逐渐成为研究热点。这类方法通过直接利用预训练的视觉编码器与大语言模型，在不更新模型参数的情况下完成任务，主要依赖提示工程与模块化设计实现跨模态协同。

在异常检测任务中，LogSAD\cite{zhang2025towards}提出了“匹配思考”（Match-of-Thought）框架，通过多模态模型捕捉视频中的逻辑规则，并结合多粒度检测与得分校准，实现稳健的零样本异常识别。在时序定位方面，TFVTG\cite{zheng2024training}利用LLM对查询语句进行语义分解，并将其与视频候选片段进行对齐，同时区分动态与静态事件以提升匹配精度。针对动作检测任务，FreeZAD\cite{han2025training}引入对比得分机制与动作性校准策略，在无监督条件下实现对未见类别的识别，并通过测试时自适应增强模型的鲁棒性。

总体来看，免训练方法通过充分挖掘预训练模型的潜在能力，在避免高昂训练成本的同时，展现出良好的泛化性能与灵活性。然而，如何进一步通过设计优秀的免训练视频理解框架以提升其在复杂场景中的稳定性与精细化能力，仍是当前研究的重要方向。本文提出的基于免训练大模型及多智能体协同的视频摘要方法正是在这一背景下展开，旨在通过纯语义驱动的方式，实现高效且通用的视频摘要生成。

\xsection{论文的主要研究内容}{Main content of the paper}


本文围绕免训练视频摘要任务展开研究，针对当前视频摘要方法对标注数据依赖强、泛化能力受限以及难以充分利用视频高层语义信息等问题，提出了一种基于大语言模型的层级化语义视频摘要框架。本文的主要研究内容如下：

1） 针对传统视频摘要方法在高层语义理解不足且依赖人工标注训练的问题，本文提出了一种基于大语言模型的层级化语义视频摘要方法。该方法首先利用核时间分割（Kernel Temporal Segmentation, KTS）算法对原始视频进行镜头划分，将连续视频流分解为具有语义独立性的场景单元，以降低长视频处理过程中的时序冗余并保留视频的结构化信息。在此基础上，本文进一步构建了“帧级描述-场景级描述-视频级描述”的多级语义表示体系：首先借助视觉-语言模型为采样帧生成细粒度文本描述，随后利用大语言模型对同一镜头内的帧描述进行聚合，生成场景级事件概括，并进一步基于场景级概括对视频进行全局摘要，得到能够表征视频主题的高层语义描述。针对视频摘要中关键片段筛选缺乏语义依据的问题，本文设计了一种基于大语言模型的语义评分策略，通过联合输入场景描述与视频级全局描述，评估各场景对整体叙事主线的贡献度，从而获得场景级重要性分数。进一步地，本文基于文本语义嵌入分别计算帧描述与所属场景描述、帧描述与视频全局描述之间的相似度，以刻画帧在局部场景和全局主题两个层面的语义贡献；随后将场景级重要性分数与帧对场景的贡献度进行耦合，并结合帧对视频主题的贡献度共同构建最终的帧级重要性序列。该方法将视频摘要任务由传统的特征选择过程转化为层级化语义理解与重要性推理过程，在无需额外监督训练的条件下，实现了兼具语义连贯性与视觉代表性的视频摘要生成。

2） 针对单一评分模型在复杂视频场景中难以同时兼顾全局主题理解与局部片段判别的问题，本文进一步提出了一种基于大语言模型的多智能体协同的视频摘要方法。该方法首先按照固定段数或固定帧窗口对长视频进行分段，并在每个片段内部均匀采样若干代表帧；随后借助视觉语言模型对同一片段内的多帧内容进行联合描述，生成面向后续推理的段级文本表示。不同于直接由单一模型对所有片段进行独立排序，本文构建了由一个规划智能体（Planner Agent）和四个专家智能体（Expert Agents）组成的协同评分框架。规划智能体首先基于全部片段描述推断视频主题，生成视频全局摘要，并动态分配叙事、视觉、情绪和信息四个维度的专家权重；在逐片段评分阶段，规划智能体结合当前片段描述与全局摘要给出全局视角下的重要性分数，而叙事专家、视觉专家、情绪专家和信息专家则分别从主线推进、视觉显著性、情绪表达和信息增量等角度对当前片段进行补充评估。
为增强长视频评分过程的时序一致性，本文进一步引入轻量级历史记忆机制，将前序片段的文本描述作为上下文信息注入后续片段的评分过程，从而缓解相邻片段之间的语义冗余与局部误判问题。在此基础上，本文将规划智能体分数与按动态权重加权后的专家分数进行融合，得到最终的段级重要性结果，并进一步将段级分数映射回原始采样帧序列，构造可直接参与标准协议评测的帧级分数并基于该分数生成视频摘要结果。该方法将视频摘要过程由单一模型的局部评分扩展为“全局规划-多专家协同-记忆增强-帧级回填”的整体推理链路，在无需额外监督训练的条件下，提高了对长视频复杂语义结构的建模能力与摘要质量。

综上所述，本文围绕免训练视频摘要任务，分别从层级化语义建模和多智能体协同推理两个角度展开研究，构建了由视频结构化表示、语义理解、重要性评分到摘要生成的完整技术链路。所提出的方法在无需额外监督训练的条件下，能够较好地兼顾视频内容的全局主题一致性与局部关键片段判别能力，为大语言模型驱动的视频摘要研究提供了一种具有可解释性与可扩展性的实现思路。

\xsection{论文的组织架构}{The organizational structure of the paper}

本文共分为五个章节，围绕免训练视频摘要方法及多智能体协同推理框架展开，具体安排如下：

第一章：绪论。介绍视频摘要任务的研究背景与意义，分析视频内容快速增长带来的信息过载问题；随后综述视频摘要及大语言模型相关研究现状，并在此基础上说明本文的研究内容、技术路线与章节安排。

第二章：相关理论基础与技术原理。主要介绍本文研究所涉及的理论基础与关键技术，包括视频特征提取与表征学习方法、大语言模型与提示工程相关机制，以及基于 Agent 的多智能体协作思想，为后续方法设计与实验分析奠定基础。

第三章：基于大语言模型的免训练视频摘要方法研究。针对传统方法在高层语义建模上的不足，提出层级化语义视频摘要框架；重点介绍基于 KTS 的视频结构化处理、帧级到场景级再到视频级的多级语义描述生成方法，以及基于场景重要性评分、帧对场景贡献和帧对视频贡献的帧级重要性建模策略；最后通过实验验证方法的有效性并分析其性能表现。

第四章：面向视频摘要的多智能体大语言模型推理框架。针对长视频摘要中单一评分模型难以兼顾全局规划与局部判别的问题，构建由 Planner Agent 与多个 Expert Agent 组成的协同推理框架；系统介绍视频分段与段级多模态描述生成、全局规划、多专家协同评分、历史记忆增强以及段级分数到帧级结果的映射与评测流程，并通过实验分析该框架在复杂视频场景下的表现。

第五章：结论与展望。对本文围绕免训练视频摘要开展的两条技术路线进行总结，归纳主要研究结论与创新点；在此基础上分析当前方法在分段策略、长程上下文建模、实验覆盖范围和系统效率等方面的不足，并对后续研究方向进行展望。